\documentclass[11pt]{article}
\usepackage[margin=1in]{geometry}
\usepackage{amsmath,amssymb,amsthm}
\usepackage{hyperref}
\hypersetup{hidelinks}

\newtheorem{theorem}{Theorem}
\newtheorem{lemma}[theorem]{Lemma}
\newtheorem{corollary}[theorem]{Corollary}
\theoremstyle{definition}
\newtheorem{definition}{Definition}

\setlength{\parskip}{4pt plus 1pt minus 1pt}
\setlength{\parindent}{0pt}

\title{Graph Structure Is Necessary for\\Information Preservation Under Bounded Context}
\author{Patrick D.\ McCarthy}
\date{}

\begin{document}
\maketitle

\begin{abstract}
Any system that (i) accumulates propositions about a world it cannot
fully observe, (ii) operates under a bounded active context --- a
hard limit on the number of propositions simultaneously available for
inference --- and (iii) must preserve information as the proposition
count grows beyond that bound, necessarily maintains a directed graph
with typed edges over its propositions. We prove this claim from two
independent arguments. The \emph{information-preservation argument}
shows that the only space-creating operation that does not lose
information is factoring --- extracting shared structure into named
nodes with typed edges --- and that factoring \emph{is} graph
construction. The \emph{retention-dynamics argument} shows that any
representation supporting sub-linear dependency queries, contextual
retrieval, and selective removal with cascade necessarily encodes
directed adjacency, which is a semantic graph regardless of substrate.
We show that continuous-weight mechanisms (e.g., transformer
self-attention) implement graph structure rather than constitute an
alternative to it. The result is substrate-independent: it applies to
biological neural networks, artificial neural networks, institutional
knowledge systems, or any other medium that stores propositions under
bounded context.
\end{abstract}

% ------------------------------------------------------------------
\section{Introduction}
\label{sec:intro}
% ------------------------------------------------------------------

Systems that accumulate knowledge face a fundamental tension. The
knowledge base grows, but the capacity to reason over it at any single
moment does not grow proportionally. Biological working memory holds
roughly $7 \pm 2$ items \cite{Miller1956}. Transformer context windows,
though large, degrade in inference quality as they fill with content
irrelevant to the current problem \cite{Liu2024}. Institutional
decision-makers cannot hold an entire organization's knowledge in a
single meeting. In each case, there is a hard or effective bound on
\emph{active context}: the number of propositions simultaneously
available for inference.

This paper asks a narrow question: \textbf{what structure must a
knowledge representation have if information is not to be lost as the
number of stored propositions exceeds the active context bound?}

The answer is a directed graph with typed edges. We prove this from two
independent arguments, each sufficient on its own:

\begin{enumerate}
\item \textbf{Information preservation} (Theorem~\ref{thm:graph-info}).
   When a bounded system must make space for new propositions, the only
   operation that creates space without losing information is
   \emph{factoring}: extracting shared structure into a named node with
   typed edges. Factoring is graph construction. Discarding without
   adjacency information is blind; lossy compression without adjacency is
   blind discarding with extra steps.

\item \textbf{Retention dynamics} (Theorem~\ref{thm:graph-dynamics}).
   Any representation that supports sub-linear dependency queries ---
   answering ``what depends on this proposition?'' in less than $O(n)$
   time --- necessarily stores directed adjacency. Selective removal
   with cascade (invalidating a proposition and propagating to its
   dependents) requires typed edges to scope correctly. These are the
   operations any system performs when maintaining knowledge under
   continuous evidence arrival.
\end{enumerate}

The result is substrate-independent. It does not prescribe an
implementation (adjacency lists, edge tables, attention masks, synaptic
connections) --- it proves that the \emph{information content} of any
adequate representation is a directed graph, regardless of how that
graph is physically realized.

\subsection{Related work}

Graph-structured knowledge representations are widely used in practice:
knowledge graphs \cite{Hogan2021}, Bayesian networks \cite{Pearl1988},
and semantic networks have a long history. The question of whether graph
structure is \emph{necessary} --- as opposed to convenient --- has
received less formal attention.

Dupoux, LeCun, and Malik \cite{DupouxLeCunMalik2026} propose a
dual-system architecture (observation + action) with meta-control
routing, motivated by cognitive science. Their framework is
architectural and prescriptive: it argues that autonomous learning
\emph{should} integrate observation, action, and meta-control. The
present work is complementary: rather than proposing what intelligent
systems should contain, we derive what they \emph{must} contain under
bounded context.

The information bottleneck principle \cite{Tishby2000} establishes that
optimal representations compress input while preserving information
relevant to output. Our information-preservation argument
(Theorem~\ref{thm:graph-info}) addresses a different question: not what
to compress, but what structure the representation must have to compress
without information loss.

Simon's bounded rationality \cite{Simon1955} establishes that cognitive
agents satisfice rather than optimize. Our bounded context constraint
($C_n$) formalizes a specific consequence: when the bound is hard,
structure is not optional but necessary.

% ------------------------------------------------------------------
\section{Formal Framework}
\label{sec:framework}
% ------------------------------------------------------------------

We introduce the minimal definitions required for the two main results.

\begin{definition}[World State Space]
\label{def:world}
Let $W$ be a measurable space of possible world states. At time $t$, the
true state $w_t \in W$ is not directly observable by any entity.
\end{definition}

\begin{definition}[Proposition and Knowledge Base]
\label{def:knowledge}
A \emph{proposition} $p$ is a claim about $W$ with an associated
likelihood function $\ell_p : W \to [0,1]$. A \emph{knowledge base} $K$
is a collection of propositions. The posterior over world states given
$K$ is:
\[
P(w \mid K) \propto P_0(w) \cdot \prod_{p \in K} \ell_p(w)
\]
where $P_0$ is a prior. The \emph{uncertainty} of $K$ about $w$ is the
surprisal $U(w, K) = -\log P(w \mid K)$ \cite{Shannon1948, Cover2006}.
\end{definition}

\begin{definition}[Bounded Active Context]
\label{def:context}
Let $C_n$ denote the \emph{active context capacity}: the maximum number
of propositions simultaneously available for inference. Let
$\mathrm{ctx}(t) \subseteq K$ be the set of propositions loaded into
active context at time $t$, with $|\mathrm{ctx}(t)| \leq C_n$.

The \emph{inference quality} for problem $f$ is the relevance density:
\[
\rho(f) = \frac{|K^{[f]}|}{|\mathrm{ctx}(t)|}
\]
where $K^{[f]} \subseteq K$ is the set of propositions relevant to $f$.
$\rho = 1$ when active context holds exactly the relevant propositions;
$\rho \to 0$ when context is dominated by irrelevant content.
\end{definition}

\begin{definition}[Growth and Curation Regimes]
\label{def:regimes}
$K$ has a maximum capacity $K_{max}$ determined by substrate constraints.
When $|K| < K_{max}$, $K$ is in the \emph{growth regime}: propositions
accumulate freely. When $|K| = K_{max}$, $K$ is in the \emph{curation
regime}: adding a new proposition requires making space.
\end{definition}

\begin{definition}[Contextual Retrieval]
\label{def:retrieval}
Given problem $f$, \emph{contextual retrieval} is the operation that
loads $K^{[f]}$ into active context. The retrieval cost is the number of
propositions examined to identify $K^{[f]}$.
\end{definition}

% ------------------------------------------------------------------
\section{Graph Necessity from Information Preservation}
\label{sec:info}
% ------------------------------------------------------------------

\begin{theorem}[Graph Necessity from Information Preservation]
\label{thm:graph-info}
Let $K$ be a knowledge base with bounded active context $C_n$
(Definition~\ref{def:context}) in the curation regime
(Definition~\ref{def:regimes}). Any representation of $K$ that
preserves information as $|K|$ grows beyond $C_n$ necessarily
maintains a directed graph with typed edges over propositions.
\end{theorem}

\begin{proof}
In the curation regime, adding any new proposition $p_{new}$ requires
making space. There are exactly three ways to make space:

\textbf{Option A: Discard an existing proposition.}
Without explicit relational structure, the entity has no mechanism to
determine which existing propositions are redundant. Two propositions
$p_1, p_2$ that share a latent factor appear as independent entries;
there is no stored evidence that they share structure. The choice of
which to discard is therefore made without information about relational
dependencies. Discarding $p_{old}$ loses the information it encodes, and
potentially loses the information that $p_{old}$ has any relation to
other propositions --- a second-order loss. This is information loss.

\textbf{Option B: Factor --- extract shared structure, reducing $|K|$
without discarding.}
If $p_1$ and $p_2$ share a latent factor $f^*$, extracting $f^*$ as a
single named node with typed edges to $p_1$ and $p_2$ reduces $|K|$ by
replacing redundant structure with a single shared reference. The
information in $p_1$ and $p_2$ is preserved; the information about their
shared grounding in $f^*$ is now explicitly stored as an edge. Space is
created without loss. Factoring also resolves the update problem: when
new evidence revises $f^*$, the single node is updated once, and the
revision propagates to all propositions that reference it via edges ---
$O(1)$ instead of $O(|K|)$.

Factoring is graph construction. Naming $f^*$, creating a node, and
storing typed edges from $p_1$ and $p_2$ to it is precisely the
operation that produces a directed graph.

\textbf{Option C: Lossy compression --- replace propositions with a
summary that occupies less space.}
This is a genuine third option, but it reduces to Option A or Option B.
Compressing $p_1$ and $p_2$ into a summary $p_s$ that discards detail is
a form of discarding: the lost detail may ground other propositions or
actions. If $p_1$ has dependents --- propositions that reference it ---
then replacing $p_1$ with $p_s$ silently invalidates those dependents
unless the entity knows what depends on the lost detail. But knowing
what depends on $p_1$ \emph{is} knowing $p_1$'s adjacency --- its
edges. Without stored adjacency, compression is blind discarding: the
entity cannot evaluate what breaks. With stored adjacency, the entity
can scope the damage, notify dependents, and compress safely --- but
storing adjacency is maintaining a graph. Compression that preserves
correctness requires the same relational structure as factoring.
Compression that ignores relational structure is Option A with extra
steps.

Any representation that makes space without factoring or correctly
scoped compression loses information; any representation that factors or
correctly scopes compression is constructing a graph.

\textbf{Consequence for alternative representations.}
A flat store cannot factor: shared structure has no representation, so
at $C_n$ it must discard.
A tree loses cross-branch propositions when one branch is pruned.
A vector store holds shared factors implicitly in embedding geometry but
cannot factor without naming --- and naming is graph construction.
Every representation that preserves information as $|K|$ grows beyond
$C_n$ must factor. Every factoring operation creates nodes with typed
edges. The graph is what factoring is.
\end{proof}

% ------------------------------------------------------------------
\section{Graph Necessity from Retention Dynamics}
\label{sec:dynamics}
% ------------------------------------------------------------------

\begin{theorem}[Graph Necessity from Retention Dynamics]
\label{thm:graph-dynamics}
Let $K$ be a knowledge base supporting continuous knowledge maintenance
under evidence arrival. The maintenance operations are:
\begin{enumerate}
\item[(i)] \textbf{Addition with provenance}: insert a proposition $p$
   together with its evidential grounds --- which existing propositions
   $p$ depends on, and by what relation.
\item[(ii)] \textbf{Contextual retrieval}: given problem $f$, retrieve
   exactly $K^{[f]}$.
\item[(iii)] \textbf{Dependency query}: given proposition $p$, determine
   which propositions are grounded through $p$.
\item[(iv)] \textbf{Selective removal with cascade}: when $p$ is
   invalidated, remove $p$ and re-evaluate every proposition whose
   grounding depends on $p$, transitively, until $K$ stabilizes.
\end{enumerate}
Any representation achieving sub-linear cost on all four operations
simultaneously necessarily encodes directed adjacency between
propositions --- and any structure encoding directed adjacency over
propositions with typed relations \emph{is} a semantic graph, regardless
of implementation.
\end{theorem}

\begin{proof}
\textbf{Part 1: Dependency queries are necessary and continuous.}
Knowledge maintenance operates continuously: every new piece of evidence
potentially shifts the status of existing propositions. When evidence
revises or invalidates a proposition $p$, the system must evaluate every
proposition whose grounding passes through $p$. This is not a batch
operation that can be deferred --- acting on a proposition whose ground
has been invalidated risks incorrect inference. The dependency query
``what is grounded through $p$?'' is therefore invoked at the frequency
of evidence arrival.

\textbf{Part 2: Lower bound without adjacency.}
Consider a representation $R$ that stores $n$ propositions without
explicit directed relationships between them (a flat store, hash table,
embedding space, or any structure that does not record which propositions
depend on which). To answer the dependency query for $p$, $R$ must
examine each of the other $n - 1$ propositions to determine whether its
grounding involves $p$. Without stored adjacency, there is no mechanism
to restrict this examination: any proposition might depend on $p$, and
the only way to determine this is to check.

The cost of a single dependency query is therefore $\Omega(n)$ for any
representation without stored adjacency. Since dependency queries are
invoked at evidence-arrival frequency (Part~1), the amortized cost of
maintenance is $\Omega(n)$ per evidence event --- growing linearly with
$|K|$.

\textbf{Part 3: Cascade amplifies the cost.}
Selective removal is not a single deletion. When $p$ is removed, each
proposition $q$ grounded through $p$ must be re-evaluated: does $q$ have
alternative grounds, or should it also be removed? If $q$ also fails,
its dependents must be re-evaluated in turn. This is a transitive
closure on the dependency relation.

With explicit adjacency, cascade follows edges: cost is $O(|D_p|)$ where
$|D_p|$ is the transitive dependency set. Without adjacency, each level
of cascade requires an $\Omega(n)$ scan. Total cost: $\Omega(n \cdot L)$
where $L$ is the cascade depth, versus $O(|D_p|)$ in the graph where
$|D_p| \ll n$ for typical propositions with bounded connectivity.

\textbf{Part 4: Any sub-linear solution encodes adjacency.}
Suppose $R$ achieves $o(n)$ cost for dependency queries. Then $R$ must
store, for each proposition, information about which other propositions
depend on it --- otherwise the $\Omega(n)$ lower bound of Part~2
applies. This stored information is a set of directed relationships: for
each proposition $p$, a record of the propositions that reference $p$ as
part of their grounding.

This is precisely a directed adjacency structure. The nodes are
propositions. The edges are directed dependency relations. Any
representation that stores this information --- regardless of whether it
is implemented as adjacency lists, edge tables, pointers, or any other
mechanism --- \emph{is} a directed graph over propositions.

\textbf{Part 5: Typed edges are necessary for correct cascade.}
A bare adjacency structure (untyped edges) is insufficient. When $p$ is
invalidated, cascade must distinguish between propositions
\emph{inferentially grounded} through $p$ (which require re-evaluation)
and propositions merely \emph{associated} with $p$ (which do not). An
untyped edge conflates these: every neighbor must be treated as
potentially dependent, degrading cascade to the full neighborhood rather
than the true dependency set.

Typed edges --- carrying the relation type (``grounds,'' ``evidences,''
``causes,'' ``supports'') --- enable cascade to follow only
dependency-bearing edges, keeping cost proportional to $|D_p|$ rather
than $|\mathrm{adj}(p)|$. Since $|\mathrm{adj}(p)|$ can be much larger
than $|D_p|$, typed edges are necessary for efficient cascade.

A directed graph with typed edges over propositions is a semantic graph.
\end{proof}

% ------------------------------------------------------------------
\section{Continuous-Weight Mechanisms Are Graph Implementations}
\label{sec:attention}
% ------------------------------------------------------------------

A natural objection: modern attention mechanisms achieve contextual
retrieval without explicit graph construction. Does this constitute a
non-graph alternative?

\begin{theorem}[Attention Mechanisms Implement Graph Structure]
\label{thm:attention}
Any continuous-weight retrieval mechanism that achieves contextual
retrieval over $n$ propositions either (a) computes directed adjacency
dynamically at $O(n^2)$ per query, or (b) stores adjacency explicitly
at $O(d)$ per query where $d$ is the local degree. In either case, the
mechanism encodes a directed graph.
\end{theorem}

\begin{proof}
The attention matrix $A_{ij} = \mathrm{softmax}(Q_i K_j^\top /
\sqrt{d_k})$ is a dense directed adjacency matrix: entry $A_{ij}$ is
the weight of the edge from proposition $i$ to proposition $j$. Full
self-attention computes the complete $n \times n$ adjacency at every
query: $O(n^2)$ per retrieval. This is the cost of \emph{not}
pre-computing the graph --- the system rediscovers dependency structure
from scratch at each step.

A pre-built graph with explicit adjacency pays $O(d)$ per retrieval,
amortizing adjacency computation over all future queries.

Sparse attention mechanisms --- learned sparsity patterns, local
windows, routing functions --- reduce the $O(n^2)$ cost by learning
which edges to keep: they are learned graph pruning. The KV cache
materializes a subgraph for reuse across queries. Flash attention
optimizes memory access patterns without changing the structure being
computed. Each optimization is a step toward the pre-built explicit
graph: reducing the cost of dynamic adjacency by progressively storing
more structure.

A fixed sparsity pattern over propositions with typed relationships is a
graph. The structural requirement is substrate-independent: whether
adjacency is stored as edge lists, learned sparse masks, pointer
structures, or synaptic weights, the information content is the same ---
which propositions are connected to which, by what relation.
\end{proof}

\begin{corollary}[Scaling Context Windows Is Self-Defeating]
\label{cor:scaling}
Growing $C_n$ directly --- the prevailing approach to scaling inference
capacity --- is self-defeating. For any specific problem $f$, $\rho$
decreases as $C_n$ fills with content not relevant to $f$
(Definition~\ref{def:context}). More context does not produce better
inference; it produces higher variance that inference must manage. The
efficient path is not to grow the context window but to grow the encoded
knowledge accessible via graph traversal --- filling the graph, not the
context window. \cite{Liu2024}
\end{corollary}

% ------------------------------------------------------------------
\section{Discussion}
\label{sec:discussion}
% ------------------------------------------------------------------

The two theorems establish graph necessity from independent premises.
Theorem~\ref{thm:graph-info} requires only the curation regime (bounded
capacity, new propositions arriving). Theorem~\ref{thm:graph-dynamics}
requires only continuous knowledge maintenance (evidence arrives,
dependencies must be traced). Any system subject to both constraints ---
which includes every physical system that accumulates knowledge --- is
doubly committed to graph structure.

\textbf{What the result does not claim.}
The theorems prove that the \emph{information content} of the
representation must be a directed graph. They do not prescribe a
specific graph formalism (property graph, RDF, Bayesian network) or
implementation (adjacency lists, sparse matrices, attention masks). The
graph is the abstract structure; the substrate is a free parameter.

The result also does not claim that all knowledge must be stored in an
\emph{explicit}, human-readable graph database. A neural network whose
learned weights encode directed dependencies between features is
maintaining a graph in the sense of Theorem~\ref{thm:graph-dynamics},
even if no symbolic node labels exist. The claim is structural, not
implementational.

\textbf{Implications for AI architecture.}
Theorem~\ref{thm:attention} reframes the relationship between
graph-based and neural approaches: they are not competing paradigms but
different implementations of the same structural necessity. Systems that
compute adjacency dynamically (full attention) pay $O(n^2)$; systems
that store it explicitly (graphs, sparse attention, KV caches) pay
$O(d)$. The trend in architecture design --- from dense attention toward
sparse, structured, and cached retrieval --- is convergence toward
explicit graph structure, driven by the same cost pressures the theorems
identify.

The result suggests that hybrid architectures combining explicit
knowledge graphs with neural retrieval are not an engineering
convenience but a response to a structural constraint: as $|K|$ grows,
systems must encode adjacency to maintain inference quality under bounded
context.

\textbf{Implications for scaling.}
Corollary~\ref{cor:scaling} challenges the prevailing approach of
growing context windows as the primary path to more capable inference.
Under bounded context, larger windows with undifferentiated content
\emph{reduce} inference quality for any specific problem. The
alternative --- encoding knowledge into a graph and retrieving the
relevant subgraph per problem --- maintains $\rho = 1$ regardless of
total knowledge size. This is the difference between scaling the window
and scaling the knowledge.

\textbf{Biological evidence.}
Biological neural networks maintain sparse, structured connectivity
rather than dense all-to-all connections. Working memory is bounded
($C_n \approx 7 \pm 2$) while long-term knowledge grows without
apparent limit \cite{Miller1956}. Retrieval is associative and
context-driven --- consistent with graph traversal from a query node
rather than exhaustive scan. The theorems suggest this architecture is
not an accident of evolutionary history but a necessary consequence of
bounded context under information-preservation pressure.

% ------------------------------------------------------------------
\section{Conclusion}
\label{sec:conclusion}
% ------------------------------------------------------------------

We have shown that graph structure over propositions is not a design
choice but a necessary consequence of two independently sufficient
constraints: information preservation under bounded capacity, and
sub-linear knowledge maintenance under continuous evidence arrival.
Continuous-weight mechanisms such as transformer self-attention implement
graph structure dynamically rather than constituting alternatives to it.
The result is substrate-independent: any system that accumulates
propositions under bounded active context and must preserve information
necessarily maintains a directed graph with typed edges, whether that
graph is realized as explicit edges, learned sparse attention, synaptic
connectivity, or any other medium.

\bibliographystyle{plain}
\begin{thebibliography}{99}

\bibitem{Cover2006}
T.~M. Cover and J.~A. Thomas.
\newblock \emph{Elements of Information Theory}.
\newblock Wiley-Interscience, 2nd edition, 2006.

\bibitem{DupouxLeCunMalik2026}
E.~Dupoux, Y.~LeCun, and J.~Malik.
\newblock Why {AI} systems don't learn and what to do about it: Lessons on
  autonomous learning from cognitive science.
\newblock \emph{arXiv preprint arXiv:2603.15381}, 2026.

\bibitem{Hogan2021}
A.~Hogan, E.~Blomqvist, M.~Cochez, C.~d'Amato, G.~de~Melo, C.~Gutierrez,
  S.~Kirrane, J.~E.~L. Gayo, R.~Navigli, S.~Neumaier, A.-C. Ngonga~Ngomo,
  A.~Polleres, S.~M. Rashid, A.~Rula, L.~Schmelzeisen, J.~Sequeda, S.~Staab,
  and A.~Zimmermann.
\newblock Knowledge graphs.
\newblock \emph{ACM Computing Surveys}, 54(4):1--37, 2021.

\bibitem{Liu2024}
N.~F. Liu, K.~Lin, J.~Hewitt, A.~Paranjape, M.~Bevilacqua, F.~Petroni, and
  P.~Liang.
\newblock Lost in the middle: How language models use long contexts.
\newblock \emph{Transactions of the Association for Computational Linguistics},
  12:157--173, 2024.

\bibitem{Miller1956}
G.~A. Miller.
\newblock The magical number seven, plus or minus two: Some limits on our
  capacity for processing information.
\newblock \emph{Psychological Review}, 63(2):81--97, 1956.

\bibitem{Pearl1988}
J.~Pearl.
\newblock \emph{Probabilistic Reasoning in Intelligent Systems: Networks of
  Plausible Inference}.
\newblock Morgan Kaufmann, 1988.

\bibitem{Shannon1948}
C.~E. Shannon.
\newblock A mathematical theory of communication.
\newblock \emph{Bell System Technical Journal}, 27(3):379--423, 1948.

\bibitem{Simon1955}
H.~A. Simon.
\newblock A behavioral model of rational choice.
\newblock \emph{The Quarterly Journal of Economics}, 69(1):99--118, 1955.

\bibitem{Tishby2000}
N.~Tishby, F.~C. Pereira, and W.~Bialek.
\newblock The information bottleneck method.
\newblock In \emph{Proceedings of the 37th Annual Allerton Conference on
  Communication, Control, and Computing}, pages 368--377, 2000.

\end{thebibliography}

\end{document}
